% Transcription — fonds Alexandre Grothendieck, shelfmark 29, batch 5 (pages 81-100).
% Established from the facsimile archives/batches/29/batch-05.pdf (a mirror of
% https://grothendieck.umontpellier.fr/29.pdf), per the skill
% transcribe-grothendieck.
%
% Fourteen of the twenty pages are transcribed. Pages 88, 90, 92, 94, 96 and 98
% are six consecutive sheets of one unrelated typescript — SGA 3, on
% diagonalisable group schemes, "Définition 1.1 ... D(M) = Hom_{S-gr}(M_S,
% G_{m,S})" through "Corollaire 3.5" and section 4 — scanned upside down
% because Grothendieck used their versos for pages 87, 89, 91, 93, 95 and 97.
% They carry nothing in his hand, so they are not transcribed and their numbers
% simply do not appear. Same treatment batch 3 gave the violet "n° 385"
% typescript.
%
% Pages 82 and 84 are versos too, but of a different kind and they are kept.
% They are two sheets of his own French typescript — page 74 and page 69 of it,
% so out of order — on the Raynaud extension, and both carry substantive
% corrections in his hand: on 84 the sentence naming "le groupe de Raynaud" is
% recast interlinearly. The criterion the earlier batches of this folder used
% is whether a verso carries something in his hand; these do. Additions made in
% hand on the typed page are marked \add{}, the typist's xxxx-ed out words and
% his own deletions \struck{}.
%
% The batch falls into four runs.
%
% Pages 81, 83 and 85 are Grothendieck's reading notes, in English, on a
% typescript whose chapters run I to X, keyed chapter by chapter and line by
% line. They continue directly from pages 78-80 of batch 4 — Murre's letter of
% 30 April 1969 sending the preliminary version of "Revêtements formels étales
% et th. de Mumford", and Murre's own two typed pages of comments on chapters I
% to X. The chapters commented on here are the same I to X, and the subject
% matter is the same, but nothing on these three sheets names the text, and
% none is supplied. They stay in English, as the correspondence of batches 1
% and 2 does.
%
% Page 86 is a sheet on inverse images of O-Modules and on infinitesimal
% neighbourhoods, in a very faint hand over heavy bleed-through; much of it is
% left \ill{}, and its crowded pencil cube is reduced to the base-change square
% that can be read.
%
% Pages 87 to 97 are the batch's substance, a six-page run Grothendieck
% paginated himself 1 to 6: the tame fundamental group of a regular formal
% scheme X along its special fibre X_o, taken as a divisor of a normal-crossings
% divisor D. The exact sequence of pages 87 and 95 puts mu^infty_{X_o,tame}
% between pi_1^t(X/D) and pi_1^t(X_o/D'_o); Lemme 1 (page 87) kills H^1 of the
% universal tame covering, Lemme 2 (page 89) identifies the class of the
% extension with the Chern class of the self-intersection of X_o, page 93 shows
% the kernel of Pic(X_m) -> Pic(X_o) uniquely divisible by n, page 95 brings in
% the Kummer theorem for X/D, and pages 95 and 97 work two examples — X_o a
% curve over an algebraically closed field, X_o = Spec Z and Spec Z - S.
%
% Pages 99 and 100 open a new question on two unpaginated sheets: a result of
% Abhyankar and a semi-global version of it (page 99), and a Kummer-theoretic
% programme for an extension 1 -> mu_n -> G -> Gamma -> 1 over a scheme carrying
% Cartier divisors D_i (page 100).
%
% The run 87-97 is written fast, with a fine nib and heavy striking, and the
% scan is 230 ppi. That is the ceiling on how much of it can be read: the
% displayed formulas come through, the connective prose often does not, and a
% good deal of pages 89, 91 and 93 is left \ill{}. That is the honest reading
% and not a shortcut.
%
% Three English words and phrases stand in the French — "normal crossings" and
% "self-intersection" on pages 87 and 95, "tame" in the subscript of mu^infty,
% "tamely ramified" on page 97 — and are left as written.
%
% The dating is the inventory's own, for the whole shelfmark.
%
% Pass: Opus 5 (claude-opus-5), 2026-08-24 — first pass, unchecked against the pages by a human.
\documentclass[11pt,a4paper]{article}
% Le préambule commun à toutes les transcriptions du fonds Grothendieck.
%
% Il tient en une idée : **l'incertitude se compose, elle ne se résout pas.**
% Une transcription qui lisse un mot douteux a détruit la seule chose pour
% laquelle on la faisait — un lecteur doit pouvoir savoir, à chaque endroit,
% ce qui était sur le feuillet et ce qui a été ajouté. Les six macros qui
% suivent sont donc l'appareil critique, et non de la décoration.
%
% Les mêmes commandes sont reconnues par `scripts/render.mjs`, qui produit la
% vue de lecture du site. Une macro ajoutée ici doit l'être aussi là-bas,
% faute de quoi le rendu échouera bruyamment — ce qui est voulu : un rendu
% silencieusement faux est pire qu'un rendu absent.

\ProvidesPackage{grothendieck}[2026/08/08 Transcriptions du fonds Grothendieck]

\usepackage{amsmath,amssymb,amsthm}
\usepackage{mathrsfs}
\usepackage{stmaryrd}
% Les diagrammes commutatifs sont partout dans ce fonds : c'est la langue
% même de la géométrie algébrique de Grothendieck, et une transcription qui
% les rendrait en prose ne transcrirait rien.
\usepackage{tikz-cd}
% Français par défaut — c'est la langue des feuillets — mais l'anglais est
% chargé aussi : la traduction et le résumé sont des documents anglais, et la
% typographie française (espaces avant les deux-points) y serait une faute.
% Ces documents basculent par \selectlanguage{english} en tête de corps.
\usepackage[english,french]{babel}
\usepackage{fontspec}
\usepackage{geometry}
\geometry{a4paper,margin=25mm,marginparwidth=28mm}
\usepackage{marginnote}
\usepackage{xcolor}
% ulem, not soul. `soul`'s \st and \ul are built on a character-by-character
% scan that XeLaTeX's font machinery breaks: the document fails with an
% undefined control sequence rather than degrading. ulem does the same two jobs
% and is Unicode-engine safe. `normalem` stops it hijacking \emph, which the
% transcriptions use for Grothendieck's own emphasis.
\usepackage[normalem]{ulem}
\usepackage{hyperref}
\hypersetup{colorlinks=true,linkcolor=black,urlcolor=black,pdfborder={0 0 0}}

% --- Métadonnées du lot -----------------------------------------------------
% Reprises telles quelles par le script de rendu, qui les affiche en tête de
% la vue de lecture. Elles fixent ce que le fichier prétend couvrir : sans
% elles, une transcription flotte sans point d'attache dans un fonds de seize
% mille feuillets.

\newcommand{\folder}[1]{\gdef\grothFolder{#1}}
\newcommand{\batch}[1]{\gdef\grothBatch{#1}}
\newcommand{\leaves}[2]{\gdef\grothFirst{#1}\gdef\grothLast{#2}}
\newcommand{\foldertitle}[1]{\gdef\grothFoldertitle{#1}}
\gdef\grothFolder{?}\gdef\grothBatch{?}\gdef\grothFirst{?}\gdef\grothLast{?}\gdef\grothFoldertitle{}

% --- Le résumé --------------------------------------------------------------
%
% Il ouvre la lecture modernisée, sous le titre « Résumé », et il est composé
% en petit. Ce n'est pas un ornement : c'est le seul passage du document qui
% ne soit pas des mathématiques, et un lecteur doit pouvoir le reconnaître
% avant de l'avoir lu — puis le sauter s'il connaît déjà le sujet, ou s'y
% arrêter s'il ne le connaît pas. Le filet qui le suit marque la reprise du
% corps.

\newenvironment{resume}
  {\small\setlength{\parskip}{0.45em}}
  {\par\medskip\hrule\medskip}

% --- Mots-clés ---------------------------------------------------------------
%
% En anglais, en clôture du résumé de la lecture modernisée : le vocabulaire
% moderne sous lequel chercher ce que les feuillets construisent. Le manifeste
% du site (`npm run manifest`) les extrait pour étiqueter la cote entière —
% cette ligne est la source unique des tags, il n'y a pas de fichier à côté.
\newcommand{\keywords}[1]{\par\smallskip\noindent
  {\footnotesize\textsc{Keywords} --- \textit{#1}}\par}

% --- L'appareil critique ----------------------------------------------------

% Le numéro de feuillet, en marge. C'est l'ancre qui permet de régler un
% désaccord sur une formule en regardant *un* feuillet plutôt qu'une chemise
% de six cents. Le site s'en sert aussi pour tourner les pages du fac-similé
% quand on fait défiler la transcription.
\newcommand{\leaf}[1]{%
  \marginnote{\footnotesize\color{gray!70}\textsf{#1}}%
  \label{leaf:#1}\ignorespaces}

% Illisible : on ne devine pas. Un mot inventé qui se lit comme les autres est
% le pire résultat possible.
\newcommand{\ill}{\textcolor{red!60!black}{[\,\ldots\,]}}

% Lecture douteuse : le mot est donné, et signalé comme incertain.
\newcommand{\uncertain}[1]{\uline{#1}}

% Ajout de l'éditeur — une lettre manquante, un mot sous-entendu.
\newcommand{\add}[1]{\textcolor{blue!50!black}{[#1]}}

% Biffé par Grothendieck lui-même. Ce qu'il a rayé fait partie du document :
% une rature dit souvent où la pensée a changé de direction.
\newcommand{\struck}[1]{\sout{#1}}

% Note du transcripteur, sur l'état matériel du feuillet.
\newcommand{\note}[1]{\par\smallskip{\small\color{gray!60!black}\textit{[#1]}}\par\smallskip}

% Note marginale de Grothendieck, distincte du corps du texte.
\newcommand{\marginal}[1]{\par\smallskip\noindent\hspace{1em}%
  {\small\color{gray!60!black}#1}\par\smallskip}

% --- Titre ------------------------------------------------------------------

\AtBeginDocument{%
  \begin{center}
    {\large\bfseries Fonds Alexandre Grothendieck --- cote n\textsuperscript{o}~\grothFolder}\\[2pt]
    {\itshape\grothFoldertitle}\\[4pt]
    {\small feuillets \grothFirst--\grothLast\ (lot \grothBatch)}\\[2pt]
    {\footnotesize\color{gray!60!black}Transcription établie d'après le fac-similé numérisé
     par l'Université de Montpellier.\\
     Les lectures incertaines sont soulignées, les illisibles notées [\,\ldots\,],
     les ajouts entre crochets.}
  \end{center}
  \vspace{1em}}

\endinput


\folder{29}
\batch{5}
\pages{81}{100}
\dating{1959-1969}
\watermark{Édition de démonstration}
\foldertitle{Groupe fondamental [Autour de SGA 4 et SGA 7] : lettres (1959, 1967, 1969), tapuscrits et copies de tapuscrit annotés (s.d.), notes manuscrites (s.d.)}

\begin{document}

\section*{Notes de lecture sur les chapitres I à X, et deux feuillets de tapuscrit annoté (pages 81 à 85)}

\page{81}
\note{Les pages 81, 83 et 85 sont des notes de lecture, en anglais, sur un
tapuscrit dont les chapitres vont de I à X ; elles sont classées par chapitre,
puis par numéro de page ou de ligne. Elles font suite aux pages 78 à 80 du
carnet 4. La colonne de gauche donne les renvois, le corps de la page les
remarques ; la disposition est rendue ici par une liste.}

\textbf{I} — 3, 7, 9, 11, 12.

restrict to \uncertain{adic} morphisms

ambiguity on point of formal scheme quite general

\textbf{I} — 18, 20, 27.

Most examples come from \ill{} gp.\ schemes over $\mathbb{Z}$

\textbf{II} — 2, 12, 16, 20, 22.

12 : unreasonable!

\textbf{III} — 1, 2, 6 (2 fois), 8, 11, 12\textsuperscript{b}, 16.

whole ch.\ \textbf{III} seems \struck{to be} used only to get 3.5.8

topological interpretation of $\pi_1^{D}$ in terms of $S$-Sch$/D$

Copy of previous notes on tame ramification?

Include variant of Lefschetz's Theorem for $\pi_1^{t}$

Include, after Mumford's theorem, application to question of finite generation
of local \struck{\ill{}} fund.\ group prime to $p$

16 : add section on case where $D$'s have equations! This should make the
situation much clearer! For instance 3.4.1, 3.4.2

\textbf{III} — 22, 25.

Appl.\ descent for tamely ramified coverings? At least one would like to have a
« champ » on Ét$_{/S}$!

\textbf{III} — 15, 26, 28.

b) Normality-assumption unnatural (formal normality, (as formal schemes\ldots))
— cf 3.6.6

Proof of 3.6.8 (?) is unbelievably long!

where is the « Abhyankar lemma »?

\page{82}
\note{Feuillet de tapuscrit, numéroté 74 par le dactylographe, portant la fin de
7.6.2 et le début de 7.6.3. Les mots surfrappés par la machine et les mots
biffés sont donnés en \struck{} ; les ajouts de la main de Grothendieck en
\add{}.}

\ldots ment constant à fibres finies et que les $_{n}G^{o}$ (pour $n \geqslant
1$) sont des schémas en groupes \struck{\ill{}} finis localement libres sur
$S$, que $H$ est fini localement libre sur $S$, et que la formation
\struck{\ill{}} de $\underline{\mathrm{Hom}}(\underline{\mathbb{Z}},
G^{o\sharp})$ commute à tout changement de base.
\note{Le caractère lu ici $\underline{\mathbb{Z}}$ est, sur la carbone, un
caractère surfrappé et souligné — un O ou un Q barré d'un trait oblique — que la
machine ne rendait manifestement pas. La lecture repose sur la phrase
elle-même : pour que la mention des $_{n}G^{o}$ finis localement libres serve à
quelque chose, il faut que $\underline{\mathrm{Hom}}$ du groupe indicé $n$ dans
$G^{o}_{n}$ soit précisément $_{n}G^{o}$, ce qui force le groupe en question à
avoir $\mathbb{Z}/n$ pour quotient. Le caractère lui-même reste douteux.}

Appliquant les mêmes réflexions aux extensions de $\underline{\mathbb{Z}}_{n}$
par $G^{o}_{n}$, on trouve que celles-ci correspondent aux torseurs sous le
Groupe $H_{n} = H \times_{S} S_{n} \simeq
\underline{\mathrm{Hom}}(\underline{\mathbb{Z}}_{n}, G^{o}_{n})$. Or, $H$ étant
fini localement libre sur $S$, on en conclut aussitôt que la
\struck{\ill{}} catégorie des torseurs sous $H$ équivaut (par le foncteur
« complétion formelle ») à la catégorie des torseurs sous $\hat{H}$, i.e.\ des
« systèmes cohérents » de torseurs $P_{n}$ sous les $H_{n}$. Cela implique donc
que la catégorie des extensions de $\underline{\mathbb{Z}}$ par $G^{o\sharp}$
équivaut à celle des \struck{extensions de} « systèmes cohérents » d'extensions
des $\underline{\mathbb{Z}}_{n}$ par les $G^{o}_{n}$, ce qui précise donc la
définition de l'extension $G^{\sharp}$ de (7.6.1).

Le schéma en groupes $G^{\sharp}$, en tant qu'extension de
$\underline{\mathbb{Z}}$ par $G^{o\sharp}$, dépend évidemment encore de façon
fonctorielle de $G$. Sa formation est compatible aux changements de base par un
morphisme de traits complets $S' \to S$.

7.6.3. Notons que les constructions précédentes gardent également un sens sur
une base plus générale qu'un \struck{\ill{}} \add{schéma local} complet ; il
suffit de partir d'un anneau noethérien adique $V$, dont la topologie est
définie par un \struck{\ill{}} idéal $\underline{m} \subset V$, et d'un schéma
en groupes $G$ sur $V$ tel que $G_{o} = G \otimes_{V} V_{o}$ (où $V_{o} = V/m$)
soit extension d'un schéma en groupes étale fini par un \struck{\ill{}} schéma
abélien par un tore. La question essentielle est celle de montrer que le schéma
formel abélien $\hat{B}$ \add{sur $V$} construit comme dans 7.1 en termes de
$G^{o}$ est algébrisable ; cela se voit par la méthode de 7.2 lorsqu'on suppose
que $S$ est normal.

\page{83}

\textbf{IV} — §1\textsuperscript{2}.

1.1.2 \uncertain{drawn} in terms of \textbf{III} p.\ 16 added.

\struck{Unnecessary}, take point of $\mathrm{Spec}(\mathcal{O}_{\mathcal{X},
s})$!

\emph{Introduce} geometric points of $\mathcal{X}$.

\textbf{VI} — 2, 7, 9, 11.

not used?! — to check!

2 : twice

7 : \struck{what} what is the meaning?

9, 11 : 6.7.1 and 6.7.2 could be made into a single example — would make a nice
example!

\textbf{VII} — 1, 2.

2 : I guess that one has equivalence of categories anyhow!!

\textbf{VIII} — 1, 2, 3/8.16.

Not too clear. Say that if \struck{$A$ a local or} \add{$S$ (any scheme)}
\struck{local} henselian ring, then to any (for instance the localization of
$\mathcal{S}$ at $s_{o}$) and \struck{$s_{o}, s_{o}^{*}$} \struck{two points}
$x, y$ two points \ill{}, then the \ill{} give one point $\bar{x}, \bar{y}$ of
$k(x)/k$, $k(y)/k$, then the « paths » of $\bar{y}$ to $\bar{x}$
\struck{$\bar{x} \otimes \bar{y}$} correspond exactly to the \struck{point}
liftings \ill{} i.e.\ to the \ill{} of $\pi_{o}(\bar{S} \otimes_{S} y) = \prod
k_{i}$, i.e.\ $k(y)$.
\note{Un petit diagramme accompagne la phrase : $\bar{S}$ = strict loc.\ of $S$
at $\bar{x}$, au-dessus de $S$, avec une flèche en pointillé venant de
$\bar{y}$.}

\textbf{IX} — 1, 4, 5, 6, 9, 10$'$, 10$''$, 12 ; 14 etc., 18, 19.

$9.3.4 \Rightarrow 9.3.5$ looks tautological

$9.3.5 \Rightarrow 9.3.2$ very cumbersome!

\marginal{yoga \ill{} Abhyankar lemma to be used}

\page{84}
\note{Feuillet de tapuscrit, numéroté 69 par le dactylographe, donc antérieur
dans le texte au feuillet de la page 82 : 7.2.1 à 7.2.5, où l'extension de
Raynaud est définie et nommée. Le passage central est repris à la main, et la
phrase biffée et sa réécriture sont données l'une et l'autre.}

Appliquons \struck{1}7.2.1 \struck{à} au groupe formel $\hat{A}$, dans le cas
favorable envisagé dans \struck{la} 7.1 où le schéma \struck{ab.} formel abélien
$\hat{B}$ est algébrisable, \add{et où on suppose de plus $T_{o}$ isotrivial}.
On trouve alors une \add{en groupes $A^{\sharp}$ sur $S$} \struck{extension
canonique du} schéma \struck{abélien $B$ dont le} \add{appelé le \emph{groupe de
Raynaud} associé au schéma en groupes $A$ sur $S_{o}$, tel que $A$} complété
formel est $B$, par le tore $T$ sur $S$ qui remonte $T_{o}$, \struck{et} que
nous appellerons l'\emph{extension de Raynaud} associée à $A$, et que nous
noterons $A^{\sharp}$ :

\[ 0 \to T \to A^{\sharp} \to B \to 0. \tag{\struck{1}7.2.2} \]

Elle est donc définie, à isomorphisme unique près, par le fait qu'elle donne
naissance à un système cohérent d'isomorphismes d'extensions

\[ A^{\sharp}_{n} \simeq A_{n} \tag{7.2.3} \]

(le deuxième membre muni de sa structure d'extension décrite dans (7.1.1)) i.e.\
qu'on a un isomorphisme de \add{schémas formels en} groupes \struck{\ill{}}

\[ (A^{\sharp})^{\wedge} \simeq \hat{A}, \tag{7.\struck{3}2.4} \]

respectant d'ailleurs les structures d'extensions. \struck{Il résulte de plus de
7.2.1.}

7.2.5. Il résulte d'autre part \struck{\ill{}} immédiatement de 7.2.1 que la
suite exacte (7.2.2) dépend fonctoriellement du schéma \add{en groupes} $A$ sur
$S$ satisfaisant aux conditions envisagées dans 7.1, et que sa formation commute
à tout changement de base par un homomorphisme $S' \to S$ associé à un
homomorphisme d'anneaux topologiques noethériens adiques $V \to V'$.

\page{85}

\textbf{X} — 1.

\ill{} works (and is useful) in quasi coherent case\ldots

State that $\mathcal{E}$ is of finite type (resp.\ finite presentation)
\add{or loc.\ free} if $\mathcal{E}_{1}, \mathcal{E}_{2}$ are.

\textbf{X} — 5.

\struck{\ill{}}
\note{Le point 5 est un croquis au crayon, entièrement barré d'une croix et
illisible ; seul « $+\,S_{3}$ » s'y lit.}

\textbf{X} — 12.

Amounts to showing: \struck{It} that \ill{} covering of $S$ is tame with r.\ to
$D_{1}, D_{2}$ iff its restriction to $S_{1}, S_{2}$ are tamely ramified —
\emph{should be trivial} with the conditions considered!

\section*{Image inverse de $\mathcal{O}$-Modules, et voisinages infinitésimaux (page 86)}

\page{86}
\note{Feuillet d'une écriture très fine, sur un papier que traverse le tapuscrit
du verso ; une grande partie n'est pas lisible à la résolution du fac-similé et
reste \ill{}.}

Voyons dans \uncertain{le} notion d'\emph{image inverse de $\mathcal{O}$-Modules}.
Un $\mathcal{O}$-Module quasi-cohérent est transformé en idem, un
\emph{Module \add{\uncertain{fini cohérent}} \uncertain{plat}} est transformé en
idem. Il faut vérifier que ceci \ill{} la notion \ill{} d'image inverse \ill{}
Modules quasi-cohérents.
\note{« dans le notion » : lecture douteuse, la syntaxe demanderait « la ».}

Supposons que \struck{si} $X' \to X \times_{S} S'$ soit \ill{} (\ill{} ou moins
ramifié), alors \ill{} le morphisme

\[ T_{o \cdot X} X' \to T_{o \cdot S} S' \]

est \ill{} \ill{} (resp.\ non ramifié), et peut être \ill{} de ce fait des
voisinages infinitésimaux de tous ordres,

\begin{tikzcd}
X \times_{S} S' \arrow[r] \arrow[d] & X \arrow[d] \\
S' \arrow[r] & S
\end{tikzcd}

\[ T_{o \cdot X} X' = T'_{o} \]

\note{La page dessine ici, au crayon et très serré, un cube dont le carré
ci-dessus n'est qu'une face : les autres sommets portent $T_{o}$, $T_{o}
\times_{S} S'$, $T'$ et deux étiquettes qui ne se lisent pas. Trois flèches
portent la mention « n.r. » (non ramifié) et un sommet est barré d'une croix. Le
cube n'est pas reproduit : l'essentiel de ses étiquettes est illisible, et le
dessiner reviendrait à affirmer ce qu'on n'a pas lu.}

dans $T'$ \ill{}. Alors $f^{*}(\mathcal{F}) = \varinjlim_{n} \xi'_{n}$ \ill{},
si $F$ est un faisceau \ill{}

\section*{$\pi_1^{t}$ d'un schéma formel régulier le long de sa fibre spéciale (pages 87 à 97)}

\page{87}
\note{Grothendieck numérote cette page 1. Les six pages de la suite — 87, 89,
91, 93, 95, 97 — portent sa pagination 1 à 6 dans l'angle supérieur droit.}

$\mathcal{X}$ schéma formel régulier \add{connexe}, $\mathcal{X}_{o} = X_{o}$
diviseur régulier, $D = \bigcup_{i = 0}^{n} D_{i}$ diviseurs \struck{\ill{}}
normal crossings deux à deux, $D_{o} = X_{o}$, \add{les $D_{i}$ réguliers
irréductibles}, $D' = \sum_{i \neq 0} D_{i}$, $D'_{o}$ induit sur $X_{o}$.
\note{La borne supérieure de la réunion définissant $D$ est peu lisible ; « $n$ »
est une lecture.}

\[ e \to I_{o} \to \pi_{1}^{t}(\mathcal{X}/D, \bar{\xi}_{1}) \to
\pi_{1}^{t}(X_{o}/D'_{o}, \bar{\xi}_{o}) \to e \]

\note{Sous la suite, deux identifications portées à la volée :
$\pi_{1}^{t}(\mathcal{X}/D)$ sous le deuxième terme, $\pi_{1}^{t}(\mathcal{U})$
et $\pi_{1}(\mathcal{X}/D', \bar{\xi}_{1})$ sous le troisième, avec
$\mathcal{U} = \mathcal{X} - D_{o}$ entouré. Un croquis en marge de gauche
dessine $D_{o} = X_{o}$ coupé par $D_{1}$ et $D_{2}$, avec $\bar{\xi}_{o}$
marqué.}

\[ \struck{\ill{}} \to \mu^{\infty}_{X_{o}, \mathrm{tame}}(\bar{\xi}_{o}) \to
I_{o} \to o \]

\marginal{$= \varprojlim_{n > o} \mu_{n}$, $n$ inversible dans
$\mathcal{O}_{X_{o}}$}

Donc

\[ \mu(\bar{\xi}_{o}) \to \pi_{1}^{t}(\mathcal{X}/D) \to
\pi_{1}^{t}(\mathcal{U}) \to e \]

\marginal{$\prod_{\ell \neq p} \mathbb{Z}_{\ell}$, $p = \mathrm{car}\, X_{o}$}

Orientations de $\pi_{1}^{t}(\mathcal{U})$, \ill{} $I_{o}$ via orientations de
$\pi_{1}^{t}(\mathcal{U})$, \struck{\ill{}} \ill{}. Extension déterminée par

\[ \alpha \struck{\ill{}} \in H^{2}(\pi_{1}^{t}(\mathcal{U}), I_{o}) \]

\emph{Lemme 1.} Soit $\tilde{\mathcal{U}}^{t}$ rev.\ \uncertain{univ.}\ de
$\mathcal{U}$ défini par $\pi_{1}^{t}$. Alors $H^{1}(\tilde{\mathcal{U}}^{t},
\mathcal{J}) = o$ si $\mathcal{J}$ première aux car.\ considérées
\uncertain{celle} de $X_{o}$ \add{\uncertain{n.\ cas} $\mathcal{J} = I$}.
\note{Deux mentions entourées accompagnent l'énoncé : « constant » et
« constant sur $\tilde{\mathcal{U}}^{t}$ ».}

\[ o \to H^{2}(\pi_{1}^{t}(\mathcal{U}), I) \xrightarrow{\;\rho\;} H^{2}(\mathcal{U}, I)
\to H^{2}(\tilde{\mathcal{U}}^{t}, I_{o})^{\pi_{1}^{t}} \to H^{3}(\pi_{1}^{t}, I) \]

\[ H^{2}(\tilde{\mathcal{U}}^{t}, I_{o})^{\pi_{1}^{t}} =
\mathrm{Hom}(\pi_{2}(\tilde{\mathcal{U}}^{t}), I)^{\pi_{1}^{t}} =
\mathrm{Hom}(\pi_{2}^{t}(\mathcal{U}), I)^{\pi_{1}^{t}} \]

\note{Le dernier terme de la suite se prolonge par une flèche verticale vers
$H^{3}(\mathcal{U}, I)$.}

Donc $\alpha$ \ill{} connu quand on connaît $\rho(\alpha)$, et l'image de plus,
dans $H^{2}(\tilde{\mathcal{U}}^{t}, I)^{\pi_{1}^{t}} \simeq
\mathrm{Hom}(\pi_{2}^{t}(\mathcal{U}), I)$ \ill{} est nulle.

\marginal{(pour $\mathcal{M}$, $H^{i}$, $i > o$)}

\marginal{true si courbe over \ill{} field}

\marginal{$H^{2}(\mathcal{U}, I) \Leftarrow H^{2}(\pi_{1}^{t}, H^{o}(\pi_{1}^{t},
I))$ \ill{} (Hochschild--Serre)}

\page{89}
\note{Grothendieck numérote cette page 2, et note en tête $i : X_{o} \to
\mathcal{X}$.}

\emph{Lemme 2.} Soit $\beta \in H^{1}(X_{o}, \mu)$ la classe de Chern
\struck{\ill{}} \uncertain{d'auto}-intersection de $X_{o}$ dans $\mathcal{X}$,
i.e.\ \struck{\ill{}}

\[ \mathcal{L}_{X_{o}/\mathcal{X}} \simeq i^{*}(\underline{L}(X_{o})). \]

Soit $\sigma : H^{1}(X, \mu) \to H^{2}(\mathcal{U}, I)$ l'homomorphisme
canonique. Alors

\[ \rho(\alpha) = \sigma(\beta). \]

Cela détermine \struck{\ill{}} l'\uncertain{extension}, $\alpha$, \ill{}, sous
réserve de \ill{} \ill{} $I$ \ill{} $\mu$.

\[ H^{2}(\mathcal{U}, \mu) \to H^{2}(\mathcal{U}, I) \to
H^{2}(\tilde{\mathcal{U}}^{t}, I_{o}) = \mathrm{Hom}(\pi_{2}^{t}(\mathcal{U}),
I_{o}) \]

\note{Sous le premier terme, $\beta|_{\mathcal{U}}$ ; à droite, « par lemme 2 ».
Au-dessus, au crayon : $H^{2}(\tilde{\mathcal{U}}, \tilde{\mu}) \simeq
\mathrm{Hom}(\pi_{1}^{t}(\ldots), \mu_{o})$.}

\ill{} sur $\rho(\alpha)$. \ill{}

\[ \pi_{2}^{t}(\mathcal{U}) = H_{2}(\tilde{\mathcal{U}}) \qquad
\struck{\ill{}} \qquad \pi_{2}^{t}(\mathcal{U}) \to H_{2}(\tilde{\mathcal{U}}^{t}) \]

\note{Un carré suit, dont les sommets lisibles sont $\beta$, $\rho(\alpha)$ et
$\mu_{o}$ ; le quatrième est biffé.}

\ldots et par lemme 1 \ill{} \ill{} dans le $\pi_{2}^{t}$ \ill{} \ill{}.

Donc \struck{\ill{}} $\mu_{o} \to I_{o}$ s'\ill{} sur $\mathrm{Im}\,
\pi_{2}^{t}(\mathcal{U})$. Je dis que \ill{} le \ill{}, i.e.\ \ill{}.

Si $I'_{o} = \mu_{o}/\mathrm{Im}\, \pi_{2}^{t}(\mathcal{U})$, on aura

\[ \mathrm{Im}\, \beta \in \ker\!\left(H^{2}(\mathcal{U}, I'_{o}) \to
H^{2}(\tilde{\mathcal{U}}^{t}, I'_{o})\right) \simeq H^{2}(\pi_{1}^{t},
I'_{o}), \]

donc $I\beta$ définit une extension \ill{}

\page{91}
\note{Grothendieck numérote cette page 3. Elle est écrite d'une plume très fine
sur un papier fortement traversé par le tapuscrit du verso, et n'est lisible
qu'en partie.}

\ldots de $\pi_{1}^{t}$ par $I'_{o}$, \ill{} \struck{\ill{}}, \ill{}
$\pi'_{1}$, \ill{} ou on a un \emph{épimorphisme}

\[ o \to \pi'_{1} \to \pi_{1}^{t}(\mathcal{X}/D) \]

et il faut prouver que c'est un isomorphisme.

Ceci revient à vérifier : ceci \ill{} \ill{} surjectif \ill{}

\[ \mu_{o} \to I_{o} \to \ill \]

\ill{} l'image \ill{} $\pi_{2}^{t}(\mathcal{U})$, \struck{\ill{}} $I_{o}$ \ill{}
de l'image de \ill{} \struck{\ill{}} \ill{} de $\mathcal{X}/X_{o}$.

L'hypothèse \uncertain{signifie} exactement que l'image de

\[ \beta \;(\in H^{2}(\mathcal{U}, \mu)) \quad\text{dans}\quad
H^{2}(\tilde{\mathcal{U}}^{t}, \mathcal{J}_{o}) \]

soit nulle ; de plus \ill{} nulle \ill{} i.e.\ \ill{} \ill{} fini \add{N.B.\
\ill{}} \ill{}

\page{93}
\note{Grothendieck numérote cette page 4. Même écriture et mêmes limites de
lecture que la page 91 ; ce qui suit est ce que les formules et les phrases
portantes laissent lire.}

\add{Ceci qui \ill{} ($\ldots I_{o} \ldots$)} \ill{} \ill{} dans le cas où la
classe $\beta$ est dans le noyau de

\[ H^{2}(\mathcal{X}, \mu_{o}) \to H^{2}(\mathcal{U}, \mathcal{J}_{o}). \]

On peut supposer plus \ill{} que $\mathcal{J}_{o}$ est d'indice $n > o$,
$\mu_{n}$ est \uncertain{opératif} sur $X_{o}$, alors $\mathcal{J}_{o} =
(\mu_{n})_{o}$, donc l'hypothèse dit que l'image de $\beta_{o} \in
\mathrm{Pic}(X_{o}) = H^{1}(X_{o}, \mathbb{G}_{m})$ dans $H^{2}(X_{o},
\mu_{n})$ soit \struck{$\beta_{o} \in \mathrm{Im}\, H^{1}$ dans $H^{2}(X_{o},
\mu_{n})$ est \ill{} $\mathrm{Im}\, H^{2}_{D_{o}}(X, \mu_{n}) = \ill$} nulle,
i.e.\ que \struck{\ill{}}

\[ \beta_{o} = n\gamma_{o} + \sum \lambda_{i} \beta_{o}^{i}, \]

où $\beta^{i}$ \ill{}, $\gamma \in \struck{\mathrm{Pic}(X_{o})}$ \ill{}, mais
\ill{}

\[ \beta \in \ill \; \mathrm{Pic}(X_{m}) \]

et on voit que le noyau de

\[ \mathrm{Pic}(X_{m}) \to \mathrm{Pic}(X_{o}) \]

est \struck{toujours} uniquement \emph{divisible} par $n$, donc il en est de
même du noyau de

\[ \mathrm{Pic}(\mathcal{X}) = \varprojlim \mathrm{Pic}(X_{m}) \to
\mathrm{Pic}(X_{o}), \]

d'où il résulte que l'on a

\[ {}_{n}\mathrm{Pic}(\mathcal{X}) \to {}_{n}\!\left(\mathrm{Pic}(\mathcal{X})/\ill\right) \]

surjectif. Donc il existe $\gamma \in \mathrm{Pic}(\mathcal{X})$ \ill{} que
$\beta \in \mathrm{Pic}(\mathcal{X})$ \ill{} $\beta_{o}$ \ill{} $n\gamma + \sum
a_{i} \beta_{i}$

\page{95}
\note{Grothendieck numérote cette page 5.}

\ldots on ait $n\gamma = \sum_{i} c_{i} \beta_{i}$.
\note{Les bornes de cette somme ne se lisent pas ; l'encadré de la marge droite
donne la même relation sous la forme $\beta^{o} = n\gamma + \sum_{1}^{r} c_{i}
\beta^{i}$, avec $\mathcal{J} \simeq \mu_{n}$.}

On utilise maintenant le théorème de Kummer pour $\mathcal{X}/D$ pour trouver un
revêtement de $\mathcal{X}/\underline{D_{o}}$ qui est \ill{} d'un \ill{} de
Galois $\mathcal{J}_{o} \simeq \mu_{n}$ \ill{} \ill{} sur $X_{o}$ \ill{}

\[ \pi_{2}^{t}(X_{o}/D'_{o}, \bar{\xi}_{o}) \to
\mu^{\infty}_{X_{o}, t}(\bar{\xi}_{o}) \to \pi_{1}^{t}(\mathcal{X}/D,
\bar{\xi}_{1}) \to \pi_{1}^{t}(X_{o}/D'_{o}, \bar{\xi}_{o}) \to o \]

$\pi_{1}^{t}(\mathcal{X}/D, \bar{\xi}_{1})$ comme extension de
$\pi_{1}^{t}(X_{o}/D'_{o}, \bar{\xi}_{o})$ est connue \uncertain{entièrement} en
termes de $X_{o}/D'_{o}$ et de la classe

\[ \alpha \in H^{2}\!\left(X_{o} \struck{- D_{o}},\;
\mu^{\infty}_{X_{o}, \mathrm{tame}}\right) \]

[self-intersection de $X_{o}$ dans $\mathcal{X}$] \ill{}

\emph{Exemple I.} $X_{o}$ courbe \ill{} sur $k$ alg.\ clos. Alors
$\pi_{2}^{t}(X_{o} - D'_{o}) = o$ \ill{} si $X_{o} \simeq \mathbb{P}^{1}_{k}$,
$D'_{o} = \emptyset$ \ill{}, on a $\pi_{2}^{t}(X_{o}, \bar{\xi}_{o}) =
\mu^{\infty}_{X_{o}, t}(\bar{\xi}_{o})$. L'\ill{} $\mathbb{Z} \to \mu^{\infty}_{t}$
\ill{} \ill{} (\ill{}). \struck{\ill{}}

1) $\left[X_{o} - D'_{o} \neq \mathbb{P}^{1}_{k}\right]$ \struck{\ill{}} :
$\pi_{1}^{t}(\mathcal{X}/D, \bar{\xi}_{1})$ est une extension \emph{triviale} de
$\pi_{1}^{t}(X_{o} - D_{o}, \bar{\xi}_{o})$ par
$\mu^{\infty}_{t}(\bar{\xi}_{o})$.

\page{97}
\note{Grothendieck numérote cette page 6.}

2°) $\left[X_{o} - D'_{o} = \mathbb{P}^{1}_{k}\right]$ i.e.\ $X_{o} =
\mathbb{P}^{1}_{k}$, $D'_{o} = \emptyset$. Alors
$\pi_{1}^{t}(\mathcal{X}/D, \bar{\xi}_{1})$ est \struck{une extension isomorphe}
à $\mu_{n}(k)$, \struck{où} $n$ est la partie première à $p = \mathrm{car}\, k$
\uncertain{du degré} de la self-intersection $X_{o} \cdot X_{o}$ dans
$\mathcal{X}$.

\marginal{[plus \ill{} dans $X_{o}$ \ill{} \ill{} \ill{}]}

\emph{Exemple II.} a) $X_{o} = \mathrm{Spec}\, \mathbb{Z}$. (Alors
$\mu^{\infty}_{t, X_{o}} = o$, donc

\[ \pi_{1}^{t}(\mathcal{X}/D, \bar{\xi}_{1}) \simeq \pi_{1}^{t}(X_{o}/D'_{o},
\bar{\xi}_{o}). \]

\struck{Le \ill{}} Si $D'_{o} = \emptyset$, on trouve $o$. Si $D'_{o} \neq
\emptyset$, la partie \uncertain{abélienne} de $\pi_{1}^{t}(X_{o}/D'_{o},
\bar{\xi}_{o})$ est calculable par la \ill{} de

\[ \left(\simeq \prod_{\ell \in D'_{o}} (\mathbb{Z}/\ell\mathbb{Z})^{*}\right), \]

\ill{} ; \ill{} doute $\pi_{1}^{t}(X_{o}/D'_{o}, \bar{\xi}_{o})$ \ill{} est
\uncertain{finie} par \uncertain{dévissage} \ill{}. Il y a des
\struck{\ill{}} \ill{} corps de nombres tamely ramified sur $\mathbb{Q}$, qui ne
sont pas \ill{} \ill{}??]

b) $X_{o} = \mathrm{Spec}\, \mathbb{Z} - S$, $S \neq \emptyset$. Alors

\[ \mu^{\infty}_{t, X_{o}} = \prod_{\ell \in S} \mu^{\infty}(\ell)_{X_{o}}. \]

Mais \ill{} $\pi_{2}^{t}(X_{o}/D'_{o})$ ???

\section*{Un résultat d'Abhyankar, et sa version semi-globale (pages 99 et 100)}

\page{99}
\note{Feuillet non paginé par Grothendieck. Le haut de la page — une phrase sur
un groupe d'ordre fini opérant sur $E$, et un carré $I_{o} \to I'_{o}$,
$I \to I'$ — est barré de plusieurs traits obliques, puis un trait horizontal
sépare ce début du programme qui suit.}

\struck{Il y a un \ill{} groupe d'ordre fini \ill{} $\pi_{1}(\mathcal{U})$ qui
opère sur $E$ via un \uncertain{p-groupe}. À fortiori,}

\begin{tikzcd}
I_{o}(\struck{\ill{}}) \arrow[r] \arrow[d] & I'_{o} \arrow[d] \\
I(\struck{\ill{}}) \arrow[r] & I'
\end{tikzcd}

a) Un résultat de Abhyankar $\oplus$

b) Version \add{semi}-globale du \struck{\ill{}} \emph{résultat
d'Abhyankar}.

$X$ schéma loc.\ \uncertain{noeth.}\ connexe, $D = \sum D_{i}$ diviseur
\uncertain{fini} \add{\ill{}} à croisements normaux (\uncertain{composantes}
$X_{i}$ \uncertain{régulières}).

On suppose $Z = \bigcap D_{i} \neq \emptyset$, et que $(X, Z)$ est
« hensélien », i.e.\ \ill{} un \uncertain{ét.} $X'$ fini sur $X$, connexe $\neq
\emptyset$, \ill{} de $Z$ dans $Z'$ \struck{$\pi_{o}(Z') \simeq
\pi_{o}(X')$}, $Z'$ \ill{} $\neq \emptyset$ \ill{}

\struck{Alors} Soit \struck{$X$} $\mathcal{U} = X - \struck{\ill{}}
\uncertain{D}$, $\xi$ un pt géom.\ de $\mathcal{U}$. Alors :

\begin{enumerate}
\item[a)] $\mathcal{U}$ connexe, et $\pi_{1}(\mathcal{U})_{t} \to \pi_{1}(X)$
surjectif.
\item[b)] Pour tout $D_{i}$, le groupe d'inertie de $\pi_{1}(\mathcal{U})_{t}$
relatif à $\pi_{1}(\mathcal{U})$ est un sous-groupe distingué
\struck{$\pi_{1}(\mathcal{U})$} commutatif $I_{i}$ de \ill{} \struck{\ill{}}, et
on a un épimorphisme canonique

\[ \struck{\ill{}} \prod_{\ell} T_{\ell}(\Omega^{*}) \to I_{i} \]

\marginal{l'inversible \struck{\ill{} $D_{i}$} (\ill{}, \ill{} $Z$)}
\item[c)] Les groupes $I_{i}$ commutent, et leur \struck{\ill{}}
\emph{produit} est le noyau de $\pi_{1}(\mathcal{U}, \xi)_{t} \to
\pi_{1}(X, \xi)$.
\end{enumerate}

a) Est immédiat.

b) L'invariance des $D_{i}$ dans un \ill{} \struck{\ill{}} \ill{} modérément
ramifié $X'$ de $X$ est \ill{} (étale local) et \emph{connexe},
\struck{donc $I$ \ill{}} irréd., donc $I_{i}$ est l'invariant. On sait bien
qu'il est $\simeq$ quotient de $\prod_{\ell} T_{\ell}(\Omega^{*})$, $\ell \neq$
car.\ \ill{} $D_{i}$. \struck{Il faut} \ill{} que $I \to$ \ill{} aux \ill{} de
$D_{i}$. C'est une histoire locale sur $D_{i}$.

c) Commutation des $I_{i}$ : \struck{\ill{}} \uncertain{c'est} local en un pt de
\ill{}

\page{100}

1) $X$ schéma, $D_{i}$ diviseurs de Cartier $\geqslant o$, $n_{i}$ entiers
$\geqslant 1$ ($1 \leqslant i \leqslant r$), $\Gamma$ schéma en groupes sur $S$,

\[ \struck{\ill{}} \; \mu_{\underline{n}} = \prod \mu_{n_{i}} \qquad
(\underline{n} = (n_{1}, \ldots, n_{r})) \quad \text{sur } X. \]

Extension de groupes sur $X$ :

\[ 1 \to \mu_{\underline{n}} \to G \to \Gamma \to 1 \]

\[ X' \to \tilde{X}_{1} \to X \]

\ill{} \ill{} sur $X$, à gr.\ $G$. On suppose \ill{} que $G$ opère sur
$\tilde{X}_{1}$, que $\mu$ y opère trivialement, que \struck{\ill{}}
$\tilde{X}_{1}$ devient \emph{torseur} \ill{} sous $\Gamma$.
\marginal{$= \mu_{\underline{n}} \tilde{X}$}

b) En \uncertain{disant} que \ill{} : opérateurs $\mu_{\tilde{X}}$, $X'$ \ill{}
sur $X$, \ill{} loc.\ $\mathcal{U}$ (sur $\tilde{X}_{1}$) de \uncertain{type}
kummerien relatif aux \ill{} \ill{} $D_{i}$ \ill{} (\ill{} \ill{} diviseurs)
(i.e.\ \ill{}) :

\[ \tilde{X}[T_{1}, \ldots, T_{r}] \big/ \left(T_{1}^{n_{1}} - a_{1}, \ldots,
T_{r}^{n_{r}} - a_{r}\right) \]

c) $\Gamma$ \struck{\ill{} les inerties des $D_{i}$ \ill{}} (locaux). On a donné
un iso

\[ \tilde{X} \times^{\Gamma} \mu \simeq \mu_{\underline{n}} \]

L'extension $G/\ill$ de $\Gamma$ par $\mu$ définit une classe $\alpha \in
H^{2}(B_{\Gamma}, \mu)$ \struck{$\ill \; H^{2}(B_{\Gamma}, \mu_{n_{i}})$}. On
\ill{} \ill{} \ill{} canonique \ill{} \ill{} $\tilde{X}$

\[ H^{2}(B_{\Gamma}, \mu) \xrightarrow{\;\varphi\;} H^{2}(X, \mu_{\underline{n}})
\simeq \prod H^{2}(X, \mu_{n_{i}}) \]

\[ H^{2}(X, \mu_{\underline{n}}) = H^{2}\!\left(B_{\Gamma}/(\tilde{X}, \Gamma),
\mu\right) \]

On veut prouver que, si $\varphi(\alpha) = (\xi_{i})_{1 \leqslant i \leqslant
r}$,

\[ \xi_{i} = c\!\left(\mathcal{O}_{X}(D_{i})\right) \pmod{n_{i}}. \]

\marginal{\ill{} \ill{} mod $n_{i}$}

2) Variantes \ill{} formelle, (ou \ill{} loc.\ annelés, \ill{})

\note{En marge de gauche, écrit dans le sens de la hauteur, un petit diagramme
reprend $B_{\Gamma}/(\tilde{X}, \Gamma) \to X$ au-dessus de $B_{\Gamma}$.}

\end{document}
